\subsection{Trabalhos Relacionados}
\label{sec:trabalhos-relacionados}

Diversos estudos têm abordado o uso de sistemas de informação no gerenciamento de transporte estudantil e universitário, com o objetivo de otimizar processos, reduzir falhas operacionais e melhorar a organização das viagens.

O trabalho PickBus, proposto por Pretto et al. \cite{pickbus2021}, apresenta uma solução multiplataforma (web e mobile) voltada ao gerenciamento de transporte fretado de passageiros, com foco em eventos e deslocamentos coletivos. O sistema permite a criação e o gerenciamento de viagens, controle de passageiros, organização de pagamentos e centralização da comunicação entre organizadores e usuários. O estudo fundamenta-se em conceitos de mobilidade urbana e empreendedorismo, demonstrando a viabilidade da solução por meio de pesquisa de mercado e testes com usuários reais. Os resultados indicam boa aceitação do sistema e evidenciam a carência de ferramentas digitais específicas para o gerenciamento de transporte fretado.

O trabalho intitulado "Sistema para Gerenciamento de Transportes", desenvolvido por Paiva \cite{anderson2020}, apresenta uma solução voltada à modernização do processo de solicitação e gerenciamento de transportes no contexto de uma instituição acadêmica. A proposta surge da necessidade de substituir métodos manuais baseados em troca de e-mails, os quais são suscetíveis a falhas de comunicação e perda de informações. O sistema permite o cadastro e a gestão de veículos, motoristas e setores, bem como a solicitação de viagens, verificação de disponibilidade, fluxo de aprovação hierárquico e alocação de recursos. Além disso, a solução contempla funcionalidades de monitoramento por meio de painéis e geração de relatórios, contribuindo para maior transparência e apoio à tomada de decisão. O estudo evidencia que a informatização do processo de gestão de transportes institucionais promove maior organização, eficiência operacional e confiabilidade das informações.

Outro trabalho relevante é o de Reis et al. \cite{reis2021transporte}, que apresenta a modelagem de um software voltado ao gerenciamento do transporte escolar público municipal. A proposta tem como foco o levantamento de requisitos e a utilização de técnicas de Engenharia de Software e UML para organizar informações relacionadas a alunos, rotas, veículos, motoristas e monitores. O estudo evidencia que a ausência de sistemas informatizados compromete a eficiência operacional e a tomada de decisão, especialmente em contextos com grande volume de usuários e necessidade de relatórios precisos. Diferentemente do presente projeto, que se concentra no transporte universitário intermunicipal e na oferta de uma API escalável para integração com diferentes aplicações cliente, o trabalho de Reis et al. (2021) enfatiza a modelagem conceitual e o contexto do ensino básico municipal, reforçando, contudo, a importância da informatização como fator essencial para a organização e confiabilidade dos serviços de transporte educacional.

O sistema proposto neste trabalho diferencia-se das soluções analisadas por abordar de forma específica o contexto do transporte universitário intermunicipal, um cenário ainda pouco explorado na literatura, sobretudo no que se refere à integração entre organização operacional e controle financeiro. Enquanto os trabalhos relacionados concentram-se no transporte escolar, no fretamento eventual para eventos ou na gestão institucional de veículos, a presente proposta foca nas particularidades do deslocamento diário e recorrente de estudantes universitários entre municípios, considerando suas demandas operacionais e administrativas.

Além da automatização do agendamento de viagens, do gerenciamento da lista de passageiros e da centralização da comunicação entre usuários e responsáveis pelo transporte, o principal diferencial do sistema está na inclusão de um módulo de controle financeiro integrado à gestão das viagens. O sistema permite a estimativa automática da arrecadação por viagem com base no número de passageiros confirmados, possibilitando ao responsável avaliar rapidamente a viabilidade financeira de cada deslocamento.

Adicionalmente, ao ser concebido como uma API escalável com modelo SaaS multi-tenant, acompanhada de uma aplicação cliente própria, o projeto amplia sua aplicabilidade, permitindo a integração com diferentes sistemas clientes, como aplicações web ou móveis, o que não é observado nos trabalhos analisados.
